\documentclass[11pt]{article}
\usepackage[margin=1in]{geometry}
\usepackage{amsmath,amssymb,amsthm}

\title{\vspace{-2cm} Zian Kang}
\date{}

\begin{document}

\maketitle

\section*{Question 1}

Statement: {Let $G$ be the graph shown, with four vertices arranged as in the figure, having $2$ parallel edges on the top, $3$ on the left, $2$ on the right, $4$ on the bottom, and one diagonal edge from the bottom-left vertex to the top-right vertex. Use the Matrix Tree Theorem to find a matrix whose determinant is $\tau(G)$. Compute $\tau(G)$.}

\medskip

\begin{enumerate}
    \item Label the vertices as follows:
\[
v_1=\text{top-left},\qquad
v_2=\text{top-right},\qquad
v_3=\text{bottom-left},\qquad
v_4=\text{bottom-right}.
\]

Then the edge multiplicities are
\[
m(v_1v_2)=2,\qquad
m(v_1v_3)=3,\qquad
m(v_2v_3)=1,\qquad
m(v_2v_4)=2,\qquad
m(v_3v_4)=4.
\]

Hence the degrees are
\[
d(v_1)=5,\qquad d(v_2)=5,\qquad d(v_3)=8,\qquad d(v_4)=6.
\]

Therefore the Laplacian matrix is
\[
L=
\begin{pmatrix}
5 & -2 & -3 & 0\\
-2 & 5 & -1 & -2\\
-3 & -1 & 8 & -4\\
0 & -2 & -4 & 6
\end{pmatrix}.
\]

\item By the Matrix Tree Theorem, $\tau(G)$ is any cofactor of $L$. Delete the fourth row and fourth column. Then
\[
L'=
\begin{pmatrix}
5 & -2 & -3\\
-2 & 5 & -1\\
-3 & -1 & 8
\end{pmatrix},
\]
so
\[
\tau(G)=\det(L').
\]

Now compute:
\[
\det(L')
=
5\begin{vmatrix}5&-1\\-1&8\end{vmatrix}
-(-2)\begin{vmatrix}-2&-1\\-3&8\end{vmatrix}
+(-3)\begin{vmatrix}-2&5\\-3&-1\end{vmatrix}.
\]

Thus
\[
\tau(G)
=
5(40-1)+2(-16-3)-3(2+15).
\]

So
\[
\tau(G)=195-38-51=106.
\]

Therefore
\[
\tau(G)=106.
\]
\end{enumerate}


\newpage

\section*{Question 2}

Statement: {Let $G$ be the $3$-regular graph with $4m$ vertices formed from $m$ pairwise disjoint kites by adding $m$ edges to link them in a ring. Prove that $\tau(G)=2m8^m$.}

\medskip

\begin{proof}
~
\begin{enumerate}
    \item Let the added edges joining consecutive kites be called link edges.

A kite is $K_4-e$, so it has $5$ edges. A spanning tree of a kite has $3$ edges. Since
\[
\binom{5}{3}=10,
\]
and exactly two $3$-edge choices form triangles, the number of spanning trees of one kite is
\[
10-2=8.
\]

\item Consider a spanning tree $T$ of the whole ring of $m$ kites:
 $T$ cannot omit two link edges, because then the graph breaks into two components. Hence $T$ omits at most one link edge. Therefore there are two cases.

\begin{enumerate}
\item $T$ uses exactly $m-1$ link edges:

Then one link edge is deleted, which breaks the ring into a chain. Inside each kite, we must choose a spanning tree. Hence the number of spanning trees of this type is
\[
m\cdot 8^m,
\]
since there are $m$ choices for the omitted link edge and $8$ choices for the spanning tree in each kite.

\item $T$ uses all $m$ link edges:

Since the whole graph has $4m$ vertices, a spanning tree has $4m-1$ edges. Using all $m$ link edges leaves
\[
(4m-1)-m=3m-1
\]
edges chosen inside the kites.

Every kite needs at least $2$ internal edges to keep its four vertices attached to the rest of the graph, and any kite with $4$ or $5$ internal edges would create a cycle. Hence exactly one kite contributes $2$ internal edges, and each of the other $m-1$ kites contributes $3$ internal edges, i.e. a spanning tree.

For the special kite, the two link vertices must not be connected internally, otherwise together with the ring of link edges we obtain a cycle. Thus we choose $2$ of the $5$ internal edges, except for the two choices that form a $2$-edge path joining the two link vertices. Hence the special kite also has
\[
\binom{5}{2}-2=8
\]
valid choices.

So the number of spanning trees of this type is also
\[
m\cdot 8^m.
\]
\end{enumerate}

\item Adding the two cases,
\[
\tau(G)=m8^m+m8^m=2m8^m.
\]
\end{enumerate}
\end{proof}

\newpage

\section*{Question 3}

Statement: {Using the Pr\"ufer code, count the following:
\begin{enumerate}
\item Count the number of trees on vertex set $[n]$ with exactly $2$ leaves.
\item Count the number of trees on vertex set $[n]$ with exactly $n-2$ leaves.
\end{enumerate}}

\medskip

\begin{enumerate}
\item In a tree on $[n]$, the number of occurrences of a label in the Pr\"ufer code is one less than its degree.

Having exactly $2$ leaves means exactly two vertices have degree $1$, and every other vertex has degree $2$. Hence the tree is a path. In the Pr\"ufer code, this means every entry is distinct, so we need an ordered list of length $n-2$ with distinct entries from $[n]$.

The number of such codes is
\[
n(n-1)\cdots 3=\frac{n!}{2}.
\]

Therefore the number of trees on $[n]$ with exactly $2$ leaves is
\[
\frac{n!}{2}.
\]

\item Having exactly $n-2$ leaves means exactly two vertices are not leaves. Hence exactly two labels appear in the Pr\"ufer code, and each must appear at least once.

Choose the two labels in
\[
\binom{n}{2}
\]
ways. Then form a code of length $n-2$ using those two labels, with both used at least once. There are
\[
2^{\,n-2}-2
\]
such sequences.

Thus the number of trees is
\[
\binom{n}{2}\bigl(2^{\,n-2}-2\bigr).
\]
\end{enumerate}

\newpage

\section*{Question 4}

Statement: {Compute $\tau(K_{2,m})$. Also compute the number of isomorphism classes of spanning trees of $K_{2,m}$.}

\begin{enumerate}
    \item Let the part of size $2$ be $\{x,y\}$ and the part of size $m$ be $\{u_1,\dots,u_m\}$. In any spanning tree $T$ of $K_{2,m}$, each $u_i$ has degree $1$ or $2$, since it is adjacent only to $x$ and $y$. Let $r$ be the number of vertices among the $u_i$ with degree $2$. Then the sum of degrees on the $\{u_i\}$ side is
\[
2r+(m-r)=m+r.
\]
But this also equals the number of edges of the tree, namely
\[
|E(T)|=(m+2)-1=m+1.
\]
Hence
\[
m+r=m+1,
\]
so
\[
r=1.
\]

Thus every spanning tree has exactly one vertex $u_i$ adjacent to both $x$ and $y$, and every other $u_j$ is attached to exactly one of $x,y$.

\item So:
\begin{enumerate}
\item choose the unique degree-$2$ vertex on the $m$-side in $m$ ways;
\item for each of the remaining $m-1$ vertices, choose whether it is attached to $x$ or to $y$.
\end{enumerate}

Therefore
\[
\tau(K_{2,m})=m\cdot 2^{m-1}.
\]

\item Now determine the isomorphism classes. Once the unique degree-$2$ vertex is fixed, the spanning tree is determined up to isomorphism by how many leaves are attached to $x$ and how many to $y$. If $a$ leaves are attached to $x$ and $b$ leaves are attached to $y$, then
\[
a+b=m-1.
\]
Swapping $x$ and $y$ gives an isomorphic tree, so only the unordered pair $\{a,b\}$ matters.

Hence the number of isomorphism classes equals the number of unordered pairs of nonnegative integers summing to $m-1$, which is
\[
\left\lfloor \frac{m+1}{2}\right\rfloor.
\]

\item Thus
\[
\tau(K_{2,m})=m2^{m-1},
\qquad
\text{and the number of isomorphism classes is }
\left\lfloor \frac{m+1}{2}\right\rfloor.
\]
\end{enumerate}

\newpage

\section*{Question 5}

Statement: {Prove or disprove: If $T$ is a minimum-weight spanning tree of a weighted graph $G$, then the $u,v$-path in $T$ is a minimum-weight $u,v$-path in $G$.}

\begin{proof}
~
\begin{enumerate}
    \item This statement is false.

\item Consider the triangle on vertices $\{u,v,w\}$ with edge weights
\[
w(uw)=1,\qquad w(wv)=1,\qquad w(uv)=\frac32.
\]

The minimum-weight spanning tree is
\[
T=\{uw,wv\},
\]
with total weight $2$.

\item However, the $u,v$-path in $T$ is
\[
u-w-v,
\]
whose total weight is
\[
1+1=2.
\]
But in $G$ there is a direct $u,v$-path consisting of the single edge $uv$, with weight
\[
\frac32<2.
\]

So the $u,v$-path in an MST need not be a minimum-weight $u,v$-path in $G$.
\end{enumerate}
\end{proof}

\newpage

\section*{Question 6}

Statement: {There are five cities in a network. The travel time for traveling directly from $i$ to $j$ is the entry $a_{ij}$ in the matrix below, with $\infty$ meaning there is no direct route. Determine the least travel time and quickest route from $i$ to $j$ for each pair $i,j$.}

\begin{enumerate}
    \item The matrix is
\[
A=
\begin{pmatrix}
0 & 10 & 20 & \infty & 17\\
7 & 0 & 5 & 22 & 33\\
14 & 13 & 0 & 15 & 27\\
30 & \infty & 17 & 0 & 10\\
\infty & 15 & 12 & 8 & 0
\end{pmatrix}.
\]

Applying Floyd's algorithm gives the all-pairs shortest-path matrix
\[
D=
\begin{pmatrix}
0 & 10 & 15 & 25 & 17\\
7 & 0 & 5 & 20 & 24\\
14 & 13 & 0 & 15 & 25\\
30 & 25 & 17 & 0 & 10\\
22 & 15 & 12 & 8 & 0
\end{pmatrix}.
\]

\item The quickest routes are:

\begin{enumerate}
\item From city $1$:
\[
1\to 2:\ 10,\quad \text{route }1-2.
\]
\[
1\to 3:\ 15,\quad \text{route }1-2-3.
\]
\[
1\to 4:\ 25,\quad \text{route }1-5-4.
\]
\[
1\to 5:\ 17,\quad \text{route }1-5.
\]

\item From city $2$:
\[
2\to 1:\ 7,\quad \text{route }2-1.
\]
\[
2\to 3:\ 5,\quad \text{route }2-3.
\]
\[
2\to 4:\ 20,\quad \text{route }2-3-4.
\]
\[
2\to 5:\ 24,\quad \text{route }2-1-5.
\]

\item From city $3$:
\[
3\to 1:\ 14,\quad \text{route }3-1.
\]
\[
3\to 2:\ 13,\quad \text{route }3-2.
\]
\[
3\to 4:\ 15,\quad \text{route }3-4.
\]
\[
3\to 5:\ 25,\quad \text{route }3-4-5.
\]

\item From city $4$:
\[
4\to 1:\ 30,\quad \text{route }4-1.
\]
\[
4\to 2:\ 25,\quad \text{route }4-5-2.
\]
\[
4\to 3:\ 17,\quad \text{route }4-3.
\]
\[
4\to 5:\ 10,\quad \text{route }4-5.
\]

\item From city $5$:
\[
5\to 1:\ 22,\quad \text{route }5-2-1.
\]
\[
5\to 2:\ 15,\quad \text{route }5-2.
\]
\[
5\to 3:\ 12,\quad \text{route }5-3.
\]
\[
5\to 4:\ 8,\quad \text{route }5-4.
\]
\end{enumerate}
\end{enumerate}

\newpage

\section*{Question 7}

Statement: {Let $T$ be a minimum-weight spanning tree in $G$, and let $T'$ be another spanning tree in $G$. Prove that $T'$ can be transformed into $T$ by a list of steps that exchange one edge of $T'$ for one edge of $T$, such that the edge set is always a spanning tree and the total weight never increases.}

\begin{proof}
~
\begin{enumerate}
    \item Induction on
\[
|E(T)\setminus E(T')|.
\]

If this number is $0$, then $T'=T$, and there is nothing to prove.

Now assume $T'\neq T$. Choose an edge
\[
e\in E(T)\setminus E(T').
\]
Add $e$ to $T'$. Since $T'$ is a tree, the graph $T'+e$ contains a unique cycle $C$.

Because $T$ is a tree, $C$ cannot lie entirely in $T$, so there exists an edge
\[
f\in E(C)\setminus E(T).
\]
Then
\[
T_1:=T'+e-f
\]
is again a spanning tree.

\item Inductive Step:
\[
w(T_1)\le w(T').
\]

Delete $e$ from $T$. This separates $T$ into two components, say with vertex sets $A$ and $B$. Since $e$ rejoins these components, it crosses the cut $(A,B)$. The cycle $C$ in $T'+e$ also contains another edge $f$ crossing this same cut. If
\[
w(f)<w(e),
\]
then replacing $e$ by $f$ in $T$ would produce a spanning tree lighter than $T$, contradicting the minimality of $T$. Hence
\[
w(e)\le w(f).
\]

Therefore
\[
w(T_1)=w(T')+w(e)-w(f)\le w(T').
\]

Also, since $e\in E(T)$ and $f\notin E(T)$, we have
\[
|E(T)\setminus E(T_1)|=|E(T)\setminus E(T')|-1.
\]
By the induction hypothesis, $T_1$ can be transformed into $T$ by such exchanges, always staying within the class of spanning trees and never increasing total weight. Prepending the exchange $T'\mapsto T_1$ completes the proof.
\end{enumerate}
\end{proof}

\newpage

\section*{Question 8}

Statement: {Minimum diameter spanning tree.
\begin{enumerate}
\item Construct a $5$-vertex example of an unweighted graph such that Dijkstra's algorithm can be run from some center and produce a spanning tree that does not have minimum diameter.
\item Construct a $4$-vertex example of a weighted graph such that Dijkstra's algorithm cannot produce an MDST when run from any vertex.
\end{enumerate}}

\medskip

\begin{enumerate}
\item Let
\[
V(G)=\{s,a,b,c,d\},
\]
and let
\[
E(G)=\{sa,sb,ac,ad,bc,bd\}.
\]

This is the graph obtained by taking the bipartite graph with parts $\{a,b\}$ and $\{c,d\}$ and then adding the edges $sa,sb$.

The vertex $s$ is a center, since
\[
d(s,a)=d(s,b)=1,\qquad d(s,c)=d(s,d)=2,
\]
so its eccentricity is $2$.

Now run Dijkstra's algorithm from $s$. Because $a$ and $b$ are both at distance $1$, and $c,d$ are both at distance $2$, tie-breaking may be done arbitrarily. One possible shortest-path tree is
\[
T=\{sa,sb,ac,bd\}.
\]
This is a shortest-path tree from $s$, but its diameter is $4$, since the path from $c$ to $d$ is
\[
c-a-s-b-d.
\]

However,
\[
T'=\{sa,sb,ac,ad\}
\]
is a spanning tree of diameter $3$.

No spanning tree can have diameter $2$, because that would have to be a star, and $G$ has no vertex adjacent to all four other vertices. Hence the minimum possible diameter is $3$, so $T'$ is an MDST, while $T$ is not.

\item Let $G=K_4$ on vertices $\{a,b,c,d\}$, with edge weights
\[
w(ab)=1,\qquad w(cd)=1,
\]
and all other edges having weight $3$.

Consider the spanning tree
\[
T=\{ab,ac,cd\}.
\]
This is the path
\[
b-a-c-d
\]
with edge weights $1,3,1$, so its diameter is
\[
1+3+1=5.
\]

Now any spanning tree on four vertices is either a star or a path. In a star, at least two incident edges have weight $3$, so the diameter is at least
\[
3+3=6.
\]
In a path, the sum of the three edge weights is at least
\[
1+1+3=5.
\]
Hence the minimum possible diameter is $5$, and $T$ is an MDST.

Now run Dijkstra's algorithm from any vertex. For example, from $a$:
\[
d(a,b)=1,\qquad d(a,c)=3,\qquad d(a,d)=3.
\]
Each of these distances is realized by the direct edge from $a$, since every alternative route to $c$ or $d$ has weight at least $4$. Thus every shortest-path tree rooted at $a$ must be the star centered at $a$:
\[
\{ab,ac,ad\},
\]
which has diameter $6$.

The same argument works from each vertex. Therefore Dijkstra's algorithm cannot produce an MDST from any starting vertex.
\end{enumerate}

\newpage

\section*{Question 9}

Statement: {Determine the minimum size of a maximal matching in the cycle $C_n$.}

\begin{proof}
~
\begin{enumerate}
    \item First, let $M$ be a maximal matching in $C_n$. Since $M$ is maximal, every edge of $C_n$ is either in $M$ or adjacent to an edge of $M$. Thus $M$ is an edge-dominating set.

Each chosen edge in a cycle can dominate at most $3$ edges of the cycle: itself and the two edges adjacent to it. Since the cycle has $n$ edges, we must have
\[
3|M|\ge n.
\]
Hence
\[
|M|\ge \left\lceil \frac{n}{3}\right\rceil.
\]

\item Now we show this bound is attained: Let the edges of $C_n$ be
\[
e_1,e_2,\dots,e_n
\]
in cyclic order.

\begin{enumerate}
\item If $n=3k$, take
\[
M=\{e_1,e_4,e_7,\dots,e_{3k-2}\}.
\]
Then $|M|=k$.

\item If $n=3k+1$, take
\[
M=\{e_1,e_4,e_7,\dots,e_{3k-2},e_{3k}\}.
\]
Then $|M|=k+1$.

\item If $n=3k+2$, take
\[
M=\{e_1,e_4,e_7,\dots,e_{3k+1}\}.
\]
Then $|M|=k+1$.
\end{enumerate}

In each case the chosen edges are pairwise disjoint, so this is a matching. Also every edge not chosen is adjacent to one of the chosen edges, so no additional edge can be added. Therefore the matching is maximal.

Thus the minimum size of a maximal matching in $C_n$ is
\[
\left\lceil \frac{n}{3}\right\rceil.
\]
\end{enumerate}
\end{proof}

\newpage

\section*{Question 10}

Statement: {Prove that $\alpha(G)\ge \dfrac{n(G)}{\Delta(G)+1}$ for any graph $G$.}

\begin{proof}
Let $S$ be a maximal independent set in $G$. Then every vertex of $G$ either lies in $S$ or is adjacent to a vertex of $S$; otherwise we could add that vertex to $S$, contradicting maximality.

Therefore
\[
V(G)\subseteq \bigcup_{v\in S} N[v],
\]
where $N[v]$ is the closed neighborhood of $v$.

For each $v\in S$,
\[
|N[v]|\le \Delta(G)+1.
\]
Hence
\[
n(G)=|V(G)|
\le \sum_{v\in S}|N[v]|
\le |S|(\Delta(G)+1).
\]
So
\[
|S|\ge \frac{n(G)}{\Delta(G)+1}.
\]

Since $S$ is an independent set and $\alpha(G)$ is the maximum size of an independent set,
\[
\alpha(G)\ge |S|\ge \frac{n(G)}{\Delta(G)+1}.
\]
\end{proof}
\end{document}